%% ============================================================
%% INTRODUCTION, OPENING -- v2 merge, drop-in for <<KEEP:intro-motivation>>.
%% Sits immediately after \section{Introduction} and immediately BEFORE the already
%% drafted framing paragraph ("Two features of the existing fidelity literature...").
%% Source: v1 lines 192-241.
%%
%% THE ONE THING THAT HAD TO CHANGE. v1's opening built, over four paragraphs, to the
%% question "does topology dominate thermo-hydraulic detail, as it does in electricity
%% transmission?" -- and the paper then answered yes. The revision's answer is that it is
%% LOSS VISIBILITY, not topology, that sets the requirement, and the title changed with
%% it. Carrying v1's build-up unaltered would set the reader up for the wrong result and
%% re-assert exactly the claim R1.1 objected to. The transmission analogy is kept, but as
%% a hypothesis the paper tests and partly overturns, which is a stronger opening than
%% one that tests and confirms.
%%
%% Editor compliance: v1 lumps removed --
%%   {Lund2014_4GDH, Werner2017, Lund2025_clean} -> Lund2014_4GDH
%%   {Frederiksen2013, Talebi2016}               -> Frederiksen2013
%%   {Frysztacki2021, FrewJacobson2016}          -> Frysztacki2021
%% ============================================================

District-heating networks serving industrial demand are central to European
decarbonisation, integrating heat pumps, combined heat and power and thermal storage in
multi-source systems~\cite{Lund2014_4GDH}, and their operational optimisation relies
increasingly on mathematical programming~\cite{PereaMoreno2026}. How much of the network
such a model represents is a design choice, and it is made along two axes: the spatial
resolution at which the network appears, from a single-bus copperplate to a pipe-by-pipe
graph; and the physical phenomena admitted, from a lossless energy balance to a model with
pressure-coupled hydraulics, temperature propagation and transport delay.

The two extremes are well populated. Copperplate models aggregate all components and
demands onto a single bus, ignoring spatial structure and thermal losses entirely --
a limitation their authors acknowledge explicitly: \citet{Wirtz2020} note that ``the
validity of the global energy balance assumption depends on the network topology, size and
the operation strategy'', without quantifying the resulting error. Full thermo-hydraulic
models resolve the network pipe by pipe and capture the thermal losses that amount to
5--15\,\% of annual heat supply~\cite{Frederiksen2013}, along with pressure-dependent
pumping, temperature propagation and the delays that follow from finite fluid velocity. A
third complication sits across both axes: extended physics introduces nonlinear and
bilinear terms incompatible with standard mixed-integer linear programming, forcing a
choice between linearisation, which preserves tractability at the cost of approximation
error, and native nonlinear formulations, which preserve the physics at the cost of solve
time.

The working assumption in practice is that finer spatial resolution and richer physics
both buy proportionate accuracy, and that the computational cost is therefore justified in
proportion. No quantitative evidence establishes that proportionality, or identifies which
axis delivers the return. The closest evidence comes from a neighbouring field: in
electricity systems, coarse spatial resolution causes capacity-expansion models to miss
intra-regional congestion and understate system cost, and the spatial aggregation choice
dominates other modelling simplifications in scenario
comparisons~\cite{Frysztacki2021}. It is natural to expect the analogue in district
heating, and that expectation is worth stating precisely because this paper finds it
holds only in part. Spatial resolution does matter -- but, as the controlled comparison
below shows, it matters as the mechanism by which a model computes its thermal losses,
not as a representation of routing. Once the losses are supplied by other means, the
spatial structure itself is worth almost nothing in dispatch cost. The distinction is
invisible to any comparison that changes both at once, which is the first of the two
limitations we address.
